% =====================================================================
\chapter{Architecture et contrat de données}
\label{ch:architecture}
% =====================================================================

\section{Le principe : une seule primitive}

L'architecture repose sur une idée simple : il existe \emph{une} transaction
produit, et toutes les autres surfaces n'en sont que des rendus ou des
répétitions.

\begin{encadre}
\textbf{La transaction unitaire} \api{POST /diagnostic/adresse} prend une
adresse en entrée et renvoie un enregistrement JSON canonique. La carte
l'affiche, le lot la répète, l'export PDF la met en page. Aucune autre surface
n'a le droit de posséder sa propre logique métier.
\end{encadre}

Cette contrainte a une conséquence forte : toute amélioration de la qualité de
la transaction unitaire bénéficie automatiquement à l'ensemble du produit. Elle
interdit aussi la classe de bugs la plus coûteuse pour un contrat de données —
deux implémentations qui divergent lentement.

\section{Vue en couches}

\begin{figure}[H]
\centering
\begin{tikzpicture}[
  scale=0.92, transform shape,
layer/.style={draw,rounded corners=3pt,minimum width=14.6cm,text width=14.6cm,
              minimum height=1.55cm,align=center,font=\small},
  arrow/.style={{Latex[length=2.4mm]}-{Latex[length=2.4mm]},thick,typhoonBlue}
]
\node[layer,fill=typhoonBlue!12,draw=typhoonBlue] (ui) at (0,0)
  {\textbf{Couche de présentation}\\[2pt]
   {\footnotesize React + TypeScript · carte Mapbox GL JS 3D · panneaux
   scénarios · export PDF · mode VFX}};

\node[layer,fill=typhoonGreen!10,draw=typhoonGreen] (svc) at (0,-2.1)
  {\textbf{Couche de service (FastAPI)}\\[2pt]
   {\footnotesize pipeline canonique · lot · cache · quotas · budget ·
   authentification · limitation de débit · rapport narratif}};

\node[layer,fill=typhoonSand,draw=typhoonGrey] (conn) at (0,-4.2)
  {\textbf{Couche connecteurs}\\[2pt]
   {\footnotesize géocodage · georisques REST · georisques WFS ·
   bdnb · flood\_alea (TRI) · hydro · meteo · panoramax · copernicus · mistral}};

\node[layer,fill=white,draw=typhoonGrey!60] (ext) at (0,-6.3)
  {\textbf{Sources externes}\\[2pt]
   {\footnotesize Géoplateforme · Géorisques · BDNB · Sandre · Open-Meteo ·
   Panoramax · Copernicus CDS}};

\draw[arrow] (ui) -- (svc);
\draw[arrow] (svc) -- (conn);
\draw[arrow] (conn) -- (ext);
\end{tikzpicture}
\caption{Architecture en quatre couches. La règle d'or : aucune couche
supérieure n'appelle une source externe directement, et aucun fait ne remonte
sans sa provenance.}
\label{fig:couches}
\end{figure}

\section{Structure du dépôt}

\begin{lstlisting}[caption={Organisation des fichiers du projet (simplifiée).},
label={lst:arbo}]
backend/
  app/api/routes/        diagnostic, flood, hydro, photo, report,
                         geocoding, health
  app/connectors/        bdnb, georisques (REST + WFS), flood_alea (TRI),
                         hydro (Sandre), meteo (Open-Meteo), copernicus,
                         geocoding (BAN), panoramax, mistral
  app/services/          canonical (pipeline unitaire), batch, budget,
                         limits, report_agent, latency_report
  app/schemas/           diagnostic_record (contrat serialise), flood,
                         hydro, photo, report, building_data
  app/core/              config, logging, api_key, rate_limit
  tests/                 24 modules de test hors-ligne

frontend/
  src/routes/            Zone (/zone), ReportPage, Dashboard, Portfolio,
                         WatchlistPage, Login, Home, Contact, Faq
  src/components/        UnifiedMap (Mapbox 3D), ScenarioPanel, RiskConsole,
                         RiskPanel, FloodConsole, LeftPanels, VfxDisclaimer
  src/zone/              moteur : floodAlea (client TRI), floodSim, damageModel,
                         impactModel, hazardSim, windSim, exposure,
                         hydroRoute, hydroLayer, pdf-export, canonicalAdapter
  src/zone/vfx/          vfxInputs (bus d'entree du mode cinematique)
  src/typhoon/           chrome applicatif, authentification, navigation

docs/workflow/           brainstorm, concept, spec, plan, scenario-plan,
                         vfx-spec, vfx-plan
constitution.md          regles non negociables
AGENTS.md                conventions et pieges, lu avant toute modification
\end{lstlisting}

\section{Le contrat canonique}

C'est l'artefact central du produit. Sa structure complète figure en
annexe~\ref{ann:contrat} ; on en présente ici la logique.

\subsection{Une adresse, treize aléas, un bâtiment}

\begin{lstlisting}[style=json,caption={Squelette du contrat canonique
(simplifié).},label={lst:contrat}]
{
  "schema_version": "1.0",
  "adresse": {
    "saisie": "Quai de la Rapee, 75012 Paris",
    "normalisee": "...",
    "citycode": "75056", "postcode": "75012", "city": "Paris",
    "lat": 48.848, "lon": 2.370,
    "geocode_score": 0.95
  },
  "aleas": [                    // 13 objets, referentiel complet
    {
      "code": "inondation",
      "libelle": "Inondation",
      "present": true,          // true | false | null (inconnu)
      "present_commune": true,
      "zonage": "Dans un perimetre PPR inondation",   // texte verbatim
      "hauteur_eau_m": 1.5,
      "per_building": {         // present uniquement si teste geometriquement
        "method": "point-in-polygon",
        "radius_m": null,
        "count": 2
      },
      "source": "georisques",
      "url_detail": "https://www.georisques.gouv.fr/...",
      "erreur": null,
      "resolution": "per-building"
    }
  ],
  "bdnb": {
    "donnees": { /* ~139 champs verbatim */ },
    "_source": {
      "provider": "BDNB", "url": "https://api.bdnb.io",
      "attribution": "Source BDNB (CSTB) - donnees a jour au <millesime>",
      "recuperee_le": "2026-08-25T10:00:00+00:00"
    }
  },
  "georisques_source": { /* provider, url, attribution, recuperee_le */ },
  "erreurs_partielles": [],
  "genere_le": "2026-08-25T10:00:01+00:00",
  "avertissement": "Ce rapport agrege les donnees publiques Georisques..."
}

// INTERDIT dans ce schema : niveau / bandes D03, tout score numerique,
// recommandations, copernicus, tout jugement composite. (FR-18, const. §2)
\end{lstlisting}

\subsection{La sémantique du champ \code{present}}

La valeur \code{present} est à trois états, et cette trichotomie porte une
partie de la valeur du produit :

\begin{table}[H]
\centering
\small
\caption{Les trois états de \code{present} et leur signification.}
\label{tab:present}
\begin{tabularx}{\textwidth}{@{}l l X@{}}
\toprule
\textbf{Valeur} & \textbf{Signification} & \textbf{Origine} \\
\midrule
\code{true} & Le point est concerné &
  Un test géométrique réussi, une proximité mesurée, ou un statut communal
  positif \\
\addlinespace
\code{false} & Le point a été testé et n'est pas concerné &
  Un test géométrique négatif, ou un statut communal négatif \\
\addlinespace
\code{null} & On ne sait pas &
  La source a échoué. Le champ \code{erreur} est rempli et une entrée est
  ajoutée à \code{erreurs\_partielles} \\
\bottomrule
\end{tabularx}
\end{table}

La règle \frref{14} interdit formellement d'inférer une valeur booléenne depuis
une source en échec. Un test de propriété le vérifie : quel que soit le mode de
panne injecté (délai d'attente, erreur 5xx, charge malformée, réponse vide), le
résultat est \code{null} avec une erreur renseignée — jamais \code{true} ou
\code{false}.

\begin{honesty}
C'est la règle qui distingue un produit de données honnête d'un produit
\og qui a l'air de marcher \fg. Un aléa absent parce que la source a échoué
ressemble exactement à un aléa absent parce qu'il n'existe pas. Sans la
trichotomie, l'utilisateur conclut \og pas de risque \fg{} là où il devrait lire
\og je ne sais pas \fg.
\end{honesty}

\subsection{La règle de sérialisation des champs optionnels}

Un détail qui a son importance pour un consommateur automatique : les champs
spécifiques à un aléa sont \emph{toujours présents} dans le JSON, à
\code{null} lorsqu'ils ne s'appliquent pas. Une clé n'est jamais absente. Cela
évite au client d'avoir à distinguer \og clé absente \fg{} de \og valeur
nulle \fg, distinction qui est une source classique de bugs d'intégration.

Une exception documentée : l'objet \code{per\_building} n'existe que si le
point a effectivement été testé géométriquement. Sa présence est elle-même
porteuse d'information.

\subsection{Le schéma comme garde-fou}

L'invariant \og pas de score \fg n'est pas une convention de style : il est
imposé par le schéma Pydantic. Un test vérifie qu'un objet contenant
temporairement un score en interne se sérialise \emph{sans} les clés
interdites. Le schéma refuse activement les champs \code{niveau}, les scores
numériques, \code{recommandations} et \code{copernicus}.

\section{Le pipeline unitaire}

\begin{figure}[H]
\centering
\begin{tikzpicture}[
  scale=0.88, transform shape,
  node distance=6mm,
  st/.style={draw=typhoonBlue,fill=typhoonBlue!6,rounded corners=2pt,
             align=center,inner sep=5pt,font=\footnotesize,text width=2.15cm},
  dec/.style={draw=typhoonAmber,fill=typhoonAmber!8,rounded corners=2pt,
             align=center,inner sep=5pt,font=\footnotesize,text width=2.15cm,
             diamond,aspect=1.9},
  arrow/.style={-{Latex[length=2mm]},thick,typhoonBlue}
]
\node[st] (addr) {Adresse saisie};
\node[st,right=of addr] (geo) {Geocodage\\{\scriptsize score $\geq$ 0,4}};
\node[st,right=of geo] (cache) {Cache ?\\{\scriptsize TTL 24 h}};
\node[st,right=of cache] (fan) {Appels paralleles\\{\scriptsize WFS + BDNB + REST}};
\node[st,right=of fan] (norm) {Normalisation\\{\scriptsize 13 aleas + resolution}};
\node[st,right=of norm] (rec) {Contrat\\{\scriptsize DiagnosticRecord}};

\foreach \a/\b in {addr/geo,geo/cache,cache/fan,fan/norm,norm/rec}
  \draw[arrow] (\a) -- (\b);

\node[dec,below=12mm of geo] (rej) {score $<$ 0,4\\422};
\node[dec,below=12mm of fan] (deg) {source en echec\\present $=$ null};
\draw[arrow,typhoonAmber] (geo) -- (rej);
\draw[arrow,typhoonAmber] (fan) -- (deg);
\draw[arrow,typhoonAmber,dashed] (deg) -- (norm);
\end{tikzpicture}
\caption{Le pipeline unitaire. Les sorties dégradées (en ambre) traversent le
même chemin que les sorties nominales : il n'existe pas de code séparé pour
les cas d'échec.}
\label{fig:pipeline}
\end{figure}

\subsection{Le géocodage d'abord, et comme point de non-retour}

Le géocodage est la première étape et constitue un point de non-retour. Sans
coordonnées résolues, il n'y a rien à joindre : une adresse introuvable ou trop
ambiguë arrête la transaction avec un HTTP 422. C'est le seul cas où le pipeline
échoue plutôt que de se dégrader, et c'est logique — la dégradation suppose
qu'il existe quelque chose à dégrader.

Le seuil de confiance est fixé à 0,4. En dessous, la réponse propose le libellé
le plus proche. Un test vérifie qu'aucun appel BDNB ou Géorisques n'a lieu dans
ce cas (compteurs de mocks).

\subsection{Collecte parallèle}

Les sources sont interrogées en parallèle via une structure de collecte
sécurisée (\code{\_safe\_call}) : chaque appel est isolé, avec son délai
d'attente et son traitement d'erreur. L'échec d'une source ne fait jamais
échouer la transaction.

\subsection{Normalisation}

La normalisation est l'étape où la valeur du produit se construit : elle
transforme les charges brutes hétérogènes en treize objets aléa homogènes,
chacun portant sa provenance et son niveau de résolution. C'est également elle
qui applique les règles forcées : RGA toujours
\code{commune-level-estimate}, sismique et radon toujours
\code{commune-level}.

\section{Les routes de l'API}

\begin{table}[H]
\centering
\footnotesize
\caption{Surface de l'API backend.}
\label{tab:routes}
\begin{tabularx}{\textwidth}{@{}l X@{}}
\toprule
\textbf{Route} & \textbf{Description} \\
\midrule
\api{POST /diagnostic/adresse} & Transaction canonique : adresse $\rightarrow$
  enregistrement complet \\
\api{POST /diagnostic/batch} & Soumission d'un lot ($\leq$ 10\,000 adresses) ;
  renvoie un identifiant \\
\api{GET /diagnostic/batch/\{id\}} & Interrogation de l'état du lot \\
\api{GET /diagnostic/adresse/rapport-pdf} & Relais vers le PDF officiel
  Géorisques (commodité souscripteur) \\
\api{GET /diagnostic/zone/building(s)} & Fiches BDNB autour du point \\
\api{GET /api/flood-alea} & Classes TRI au point (WFS, point-in-polygon) \\
\api{GET /api/hydro}, \api{/api/hydro/extend} & Tracé hydrographique (Sandre) \\
\api{GET /api/meteo} & Pluie horaire + débit GloFAS (Open-Meteo) \\
\api{GET /api/photo} & Photographie Panoramax et crédit \\
\api{GET /api/report}, \api{/api/report/stream} & Rapport narratif (Mistral,
  repli sur gabarit si la clé est absente) \\
\api{GET /api/geocode/search} & Autocomplétion d'adresse (BAN) \\
\api{GET /health}, \api{/health/detailed} & Sonde de vie (non authentifiée) \\
\bottomrule
\end{tabularx}
\end{table}

\section{Sécurité}

\subsection{Authentification par clé}

L'authentification est volontairement simple : une clé statique transmise dans
l'en-tête \code{X-API-Key}, validée contre une liste configurée par variable
d'environnement. Une clé absente produit un 401 ; une clé invalide, un 403. Les
routes de santé sont exemptées.

Un choix explicite : \textbf{les clés ne sont jamais journalisées}. Un test
capture les journaux et vérifie que le matériel de clé est absent. Les adresses,
en revanche, peuvent apparaître dans les journaux applicatifs — elles
constituent l'interface de requête, et les tracer est nécessaire pour le
diagnostic.

\subsection{Défenses en profondeur}

\begin{itemize}[leftmargin=1.4em]
  \item Tous les appels externes portent un délai d'attente explicite. Aucune
        attente non bornée n'existe dans le code.
  \item Une limitation de débit entrante protège l'API \typhoon{} elle-même,
        distincte de la conformité aux quotas amont.
  \item Aucun secret dans le code : toutes les informations d'authentification
        passent par des variables d'environnement, et le fichier
        \fichier{.env.example} documente chaque variable.
\end{itemize}

\section{Le mode dégradé}

Le traitement des pannes est conçu comme une fonctionnalité de premier ordre.
Le tableau~\ref{tab:erreurs} résume le comportement attendu dans chaque
situation.

\begin{table}[H]
\centering
\small
\caption{Matrice des erreurs : ce que le client observe dans chaque situation.}
\label{tab:erreurs}
\begin{tabularx}{\textwidth}{@{}l X@{}}
\toprule
\textbf{Condition} & \textbf{Comportement observable} \\
\midrule
Adresse introuvable & 422 \code{adresse\_non\_trouvee}, avant tout appel source \\
Score de géocodage $<$ 0,4 & 422 \code{adresse\_ambigue} avec libellé proposé \\
Géocodeur indisponible & 429 honoré une fois via \code{retry-after}, puis 502
  \code{geocodage\_indisponible} \\
Une couche WFS en échec & Cet aléa : \code{present=null}, \code{erreur}
  renseignée, entrée dans \code{erreurs\_partielles} — réponse 200 partielle \\
Géorisques entièrement indisponible & Les faits communaux disponibles sont
  conservés ; les faits géométriques passent à \code{null} \\
BDNB indisponible & \code{bdnb=null} + entrée dans \code{erreurs\_partielles} \\
Clé absente ou invalide & 401 / 403, avant tout traitement \\
Échec d'un élément de lot & L'élément est signalé individuellement, le lot
  continue \\
Exception inattendue & Journalisée avec trace, 502 structuré — jamais de 500 HTML \\
\bottomrule
\end{tabularx}
\end{table}

La règle de conception sous-jacente : \emph{le mode dégradé est testé avec la
même rigueur que le chemin nominal}. Une suite de tests dite
\og chaos \fg{} éteint chaque source indépendamment, dans chacun de ses modes
de défaillance, et vérifie que la réponse reste un 200 partiel structuré.
